\subsection{比例零軸の端点は固定 \texorpdfstring{$s$}{s} だけで有限になる}
\label{nhpe25:section}
比例零軸 $x=y=0,z\ge1,s<t,0\le n\le z$ の端点に対する必要条件は
\[
 n=0:\quad s\ge3,\quad s<b<3s/2,\quad 1\le z<s/2,
\]
\[
 n=z:\quad s\ge2,\quad 1\le z\le s,\quad
 1\le e=C_z/(b+1)\le s-1.
\]
算術稿の定理 \texttt{hpe25:main} に完全証明を置く．
$n=0$ では各 $b$ について元の pair 四次式を $z$ の真の二次式として解くので，
全原式検査前の候補は高々 $2\lfloor(s-1)/2\rfloor$ 個である．
$n=z$ では $e$ の四次式が正の半直線で厳密増加し，固定 $(s,z)$ に正整数根は高々一つである．
$s$ 自体に一様上界はない．両端点全体の不存在も内部 $0<n<z$ も未解決である．

\paragraph{今後使う正整数商の座標．}
\[
 C=C_n=n((sz)^2-3sz+3)+(s^2-1)z-3s+1,\qquad e=C/(b+1)>0
\]
と置くと
\[
 t=s+n(b+2sz-2)+e,\qquad S=s+n(sz-2)+e.
\]
一般の内部にも使える変換として，$b=C/e-1$ を元残余へ代入すると
\[
 e^2\mathfrak r=e^4+A e^3+B e^2+CH e-DC^2,
 \qquad D=1+n(z-n)>0,
\]
\begin{align*}
 A={}&3nsz-6n+3s+2z-1,\\
 B={}&4n^2s^2z^2-15n^2sz+15n^2-ns^2z^3+7ns^2z\\
 &+7nsz^2-2nsz-15ns-11nz+5n-s^2z^2+3s^2\\
 &+7sz-2s+z^2-z-2,\\
 H={}&3n^2sz-6n^2-2nsz^2+3ns+5nz-n-2sz+3.
\end{align*}
これらの全係数は元残余との恒等式として検算済みである．
内部では $B,H$ の符号が一定でないので，$n=z$ の単調性をそのまま移さない．
次の候補は $k=z-n$ を固定した境界近傍，または整数根を隣接整数間に挟む議論である．

\paragraph{検算を利用するときの範囲．}
\nolinkurl{verify_high_d_axis_endpoints_v25.py} の \texttt{--s-cap 400} は，
$1\le s\le400$ の全端点候補をこの有限化で覆う．
$n=0$ の $b$ は $39{,}800$ 個，$n=z$ の $(z,e)$ は $21{,}333{,}200$ 組である．
後者は単調性に基づく二分探索により，全格子を一つずつ走査せず完全に判定する．
両端点とも整数多項式根と模型は零個だった．これはこの $s$ 範囲だけの結果であり，
無上界の不存在へ外挿しない．旧 $144$ 低行検算の繰り返しでもない．

元残余を保持すべき具体例は
$(s,b,z,n,t,e)=(3,2,1,1,10,1)$ である．差方程式，整数商，尺度と全符号は通るが，
$\mathfrak r=81$ なので模型ではない．合同条件と符号だけで候補を数える誤りを防ぐために残す．
